\documentclass[a4paper,11pt]{article}
\usepackage[utf8]{inputenc}
\usepackage[T1]{fontenc}
\usepackage[ngerman]{babel}
\usepackage{amsmath,amssymb}
\usepackage{xcolor}
\usepackage{colortbl}
\usepackage{array}
\usepackage{tikz}
\usepackage{enumitem}
\usepackage[margin=2cm]{geometry}
\definecolor{bgdark}{HTML}{1E1E1E}
\definecolor{fgtext}{HTML}{E6E6E6}
\definecolor{gridcol}{HTML}{4A4A4A}
\definecolor{accent}{HTML}{6CB6FF}
\definecolor{loescol}{HTML}{7FD18C}
\pagecolor{bgdark}
\color{fgtext}
\arrayrulecolor{fgtext}
\setlength{\parindent}{0pt}
\newcommand{\karo}[1]{\par\noindent\begin{tikzpicture}\draw[step=5mm,gridcol,very thin](0,0) grid (\linewidth,#1);\end{tikzpicture}\par\vspace{2mm}}
\newcommand{\aufgabe}[2]{\vspace{4mm}\par\noindent{\large\color{accent}\textbf{Aufgabe #1}}\hfill{\color{gridcol}\normalsize #2}\par\vspace{1mm}}
\newcommand{\titel}[1]{\begin{center}{\LARGE\color{accent}\textbf{Optimierungsverfahren}}\\[3pt]{\large #1}\end{center}\vspace{1mm}\noindent\textbf{Name:}\ \underline{\hspace{6cm}}\hfill\textbf{Datum:}\ \underline{\hspace{4cm}}\\[3pt]{\color{gridcol}\rule{\linewidth}{0.4pt}}\vspace{2mm}}

\begin{document}
\titel{Übungsaufgaben --- Blatt 01: Grundlagen \& Lineare Optimierung}

\aufgabe{1 --- Klassifikation}{leicht}
Klassifizieren Sie die folgenden Optimierungsprobleme jeweils nach: (i) \emph{linear} oder \emph{nichtlinear},
(ii) Art der Nebenbedingungen (ohne / Grenzen / Gleichung / Ungleichung), (iii) Variablen ($\mathbb{R}^n$ oder $\mathbb{Z}^n$),
(iv) ein- oder mehrkriteriell.
\begin{enumerate}[label=(\alph*)]
  \item $\min\limits_{\vec x}\ x_1^2+x_2^2 \quad\text{u.d.N.}\quad x_1+x_2=1$
  \item $\max\limits_{\vec x}\ 3x_1+5x_2 \quad\text{u.d.N.}\quad x_1+x_2\le 10,\ \ x_1,x_2\ge 0$
  \item $\min\limits_{\vec x}\ \big(f_1=x_1,\ f_2=x_2\big) \quad\text{u.d.N.}\quad x_1^2+x_2^2\ge 1,\ \ x_1,x_2\in\mathbb{Z}$
  \item $\max\limits_{\vec x}\ x_1\cdot x_2 \quad\text{u.d.N.}\quad 0\le x_1\le 4,\ \ 0\le x_2\le 4$
\end{enumerate}
\karo{5cm}

\aufgabe{2 --- Modellierung}{leicht--mittel}
Eine Schreinerei stellt Tische und Stühle her. Ein Tisch bringt $40$ Euro Gewinn, ein Stuhl $25$ Euro.
Ein Tisch benötigt $4$ Arbeitsstunden und $5$ Bretter, ein Stuhl $3$ Arbeitsstunden und $2$ Bretter.
Pro Woche stehen $240$ Arbeitsstunden und $200$ Bretter zur Verfügung.\\[2pt]
Definieren Sie die Entscheidungsvariablen (mit Einheit) und stellen Sie das vollständige Optimierungsproblem
(Zielfunktion + Nebenbedingungen) auf. Ist es linear? (Eine Lösung ist \emph{nicht} nötig.)
\karo{6cm}

\aufgabe{3 --- Standardform}{mittel}
Überführen Sie das folgende Problem in die \textbf{Standardform} der linearen Optimierung
(nur Minimierung, nur Gleichungs-NB, alle Variablen $\ge 0$, rechte Seiten $\ge 0$):
\[
\max_{\vec x}\ 5x_1-2x_2+x_3 \quad\text{u.d.N.}\quad
\begin{aligned}
 x_1+x_2 &\le 6\\
 2x_1-x_3 &\ge 3\\
 x_1+x_2+x_3 &= 5
\end{aligned}
\qquad x_1\ge 0,\ x_2\ge 0,\ x_3\ \text{frei}.
\]
\karo{6.5cm}

\aufgabe{4 --- Grafische Lösung}{mittel}
Lösen Sie das folgende lineare Problem \textbf{grafisch}. Zeichnen Sie den zulässigen Bereich, verschieben
Sie die Isolinie der Zielfunktion und lesen Sie das Optimum $\vec x^\ast$ sowie den Zielfunktionswert $f(\vec x^\ast)$ ab.
\[
\max_{\vec x}\ f(\vec x)=2x_1+3x_2 \quad\text{u.d.N.}\quad x_1+x_2\le 4,\ \ x_1+3x_2\le 6,\ \ x_1,x_2\ge 0.
\]
\karo{8cm}

\aufgabe{5 --- Eckpunkte rechnerisch}{mittel}
Bestimmen Sie für das folgende Problem \textbf{alle zulässigen Eckpunkte rechnerisch}
(Schnittpunkte je zweier Begrenzungsgeraden, inkl. der Achsen), prüfen Sie die Zulässigkeit
und geben Sie das Optimum an.
\[
\min_{\vec x}\ f(\vec x)=3x_1+2x_2 \quad\text{u.d.N.}\quad x_1+x_2\ge 4,\ \ x_1+3x_2\ge 6,\ \ x_1,x_2\ge 0.
\]
\karo{7.5cm}

\aufgabe{6 --- Slackvariablen}{mittel}
Gegeben sei das lineare Problem
\[
\max_{\vec x}\ 2x_1+3x_2 \quad\text{u.d.N.}\quad
g_1:\ x_1+x_2\le 4,\quad g_2:\ x_1+3x_2\le 6,\quad g_3:\ x_2\le 3,\quad x_1,x_2\ge 0,
\]
dessen Optimum bei $\vec x^\ast=(3,1)$ liegt. Führen Sie für $g_1,g_2,g_3$ je eine Slackvariable ein.
Welche Nebenbedingungen sind im Optimum \emph{aktiv}, welche \emph{inaktiv}? Geben Sie die Werte $s_1,s_2,s_3$ an.
\karo{5cm}

\aufgabe{7 --- Ganzzahlige Optimierung}{mittel--schwer}
Betrachten Sie das ganzzahlige Problem
\[
\max_{\vec x}\ f(\vec x)=x_1+2x_2 \quad\text{u.d.N.}\quad x_1+x_2\le 5,\ \ x_2\le 2{,}5,\ \ x_1,x_2\ge 0,\ \ x_1,x_2\in\mathbb{Z}.
\]
\begin{enumerate}[label=(\alph*)]
  \item Lösen Sie zunächst die \emph{kontinuierliche} Relaxation (ohne Ganzzahligkeit).
  \item Runden Sie die Variablen jeweils auf die nächste zulässige ganze Zahl \emph{ab}. Welcher Wert ergibt sich?
  \item Bestimmen Sie das echte ganzzahlige Optimum. Was zeigt der Vergleich zu (b)?
\end{enumerate}
\karo{6.5cm}

\aufgabe{8 --- Begriffe}{leicht}
Beantworten Sie kurz:
\begin{enumerate}[label=(\alph*)]
  \item Was bedeutet \emph{unimodal} vs.\ \emph{multimodal}? Geben Sie je ein Beispiel einer Funktion.
  \item Warum genügt bei einem \emph{konvexen} Minimierungsproblem ein \emph{lokaler} Optimierer?
  \item Ist $f(x)=\sin(x)$ auf $[0,4\pi]$ unimodal oder multimodal? Begründen Sie.
\end{enumerate}
\karo{5.5cm}

\end{document}
